% !TeX root = ../main.tex

% ==================================================
% 1.1 Introduction
% ==================================================

\subsection{Python 실행 환경과 외부 시스템 개요}

\subsubsection{\zsh{sys}는 파일 경로용 패키지가 아니다}

\begin{dhknote}
\begin{dhksmry}
sys $\neq$ 시스템의 파일이나 디렉터리를 다루는 패키지

sys = 현재 실행 중인 Python interpreter와 
      그 실행 환경에 접근하는 모듈
\end{dhksmry}
\end{dhknote}

\begin{dhkcode}
import sys

print(sys.version) # 현재 어떤 Python인가? 
print(sys.executable) # 현재 어떤 Python 실행 파일을 쓰는가?
print(sys.argv) # 이 프로그램을 실행할 때 무엇을 전달받았는가?
print(sys.path) # import할 모듈을 어디서 찾는가?
\end{dhkcode}

위 4개의 print문의 인수는 전부 파일 조작 기능이 아니다.

\begin{dhkcode}
from pathlib import Path # 파일 경로 다루는 패키지
import os # 운영체제 기능을 다루는 패키지
\end{dhkcode}

그러므로 아래와 같이 생각하자.

\begin{dhkdir}
sys
>Python interpreter
>실행 인수
>import 검색 경로
>stdin / stdout / stderr
>프로그램 종료 등
\end{dhkdir}

\begin{dhkdir}
os
>운영 체제
>환경 변수
>프로세스 관련 기능
>일부 파일 시스템 기능
\end{dhkdir}

\begin{dhkdir}
pathlib
>파일/디렉터리 경로
\end{dhkdir}

\subsubsection{그런데 왜 sys를 배우다 보면 경로가 나오는가?}

대표적인 원인은 \texttt{sys.path}이다.

\begin{dhknote}
\begin{dhksmry}[centered]
sys.path $\neq$ 내 컴퓨터의 파일 경로를 탐색하는 기능

sys.path = \textbf{Python이 import 할 모듈을 찾는 위치들의 목록}
\end{dhksmry}
\end{dhknote}

예를 들어:

\begin{dhkcode}
import numpy
\end{dhkcode}

라는 코드를 실행했다고 하면 Python 입장에서는 문제가 발생한다.

\begin{dhksmry}[centered]
numpy가 어디 있는데?
\end{dhksmry}

컴퓨터에는 수십만 개의 파일이 있을 수 있다.

Python이 C 드라이브나 SSD 전부를 뒤지는 것이 아니다.

Python은 정해진 검색 위치를 가지고 있고 그것을 확인하는 게:

\begin{dhkcode}
import sys

print(sys.path)
\end{dhkcode}

위 코드의 결과는 대략 아래와 같다.

\begin{dhkzsh}
[
'/Users/user/project',
'/Users/user/anaconda3/envs/transformer/lib/python310.zip',
'/Users/user/anaconda3/envs/transformer/lib/python3.10',
...
'/Users/user/anaconda3/envs/transformer/lib/python3.10/site-packages'
]
\end{dhkzsh}

여기에서 \zsh{site-packages}는 \zsh{pip install transformers}같은 걸 하면 설치된 Python 패키지가 들어가는 대표적인 위치다.

그렇기에:

\begin{dhkcode}
import transformers
\end{dhkcode}

가 가능한 것이다.

\subsubsection{이것을 conda 가상환경과 연결해보자}

윈도우에서 Conda 환경이:

\begin{dhkzsh}
C:\Users\user\anaconda3\envs\transformer\
\end{dhkzsh}

에 있을 때 터미널에서 conda 가상환경을 활성화하고 파이썬에서:

\begin{dhkcode}
import sys

print(sys.executable)
\end{dhkcode}

을 하면 결과는

\begin{dhkzsh}
C:\Users\user\anaconda3\envs\transformer\python.exe
\end{dhkzsh}

가 나온다.

즉,

\begin{dhksmry}[centered]
Conda environment
\da
transformer
\da
그 환경의 Python interpreter
\da
\zsh{C:\Users\user\anaconda3\envs\transformer\python.exe}
\end{dhksmry}

가 연결된다.

\subsubsection{파일 경로는 또 다른 문제다}

예를 들어 프로젝트가:

\begin{dhkdir}[in=8]
transformer/
>notebooks/
>>attention.ipynb
>data/
>>emoitons.csv
>model/
>train.py
\end{dhkdir}

라고 할 때, \zsh{train.py}에서 데이터를 읽고 싶다면:

\begin{dhkcode}
import pandas as pd

df = pd.read_csv(
    'data/emotions.csv'
)
\end{dhkcode}

여기서 \zsh{'data/emotions.csv'}는 \zsh{sys.path}와 직접적인 관계가 없다.

이 것은 \textbf{파일 시스템 경로}다.

그래서 중요한 개념이 등장한다.

\begin{dhknote}
\begin{dhksmry}[centered]
현재 작업 디렉터리

Current Working Directory

CWD
\end{dhksmry}
\end{dhknote}

\subsubsection{그래서 우리가 제일 먼저 깊게 배워야 하는 것은 \zsh{sys}가 아니다}

추천 순서는:

\begin{dhksmry}[centered]
파일 시스템
\da
경로
\da
절대 경로 / 상대 경로
\da
현재 작업 디렉터리
\da
pathlib
\da
Python 프로그램 실행
\da
sys.argv
\da
sys.executable
\da
import
\da
sys.path
\end{dhksmry}

\subsubsection{최종적으로는 이런 코드를 이해하는 것이 목표다}

예를 들어 몇 단계 뒤에는 터미널에서:

\begin{dhkzsh}
python train.py \
    --data ./data/emotions.csv \
    --model distilbert-base-uncased \
    --epochs 5 \
    --output ./models/emotion-model
\end{dhkzsh}

를 실행할 때:

\begin{dhkbranch}[in=5]
node main :: Shell
arrow main :: python train.py --data ...

node main :: Python interpreter
arrow main :: sys.argv

node main :: train.py

branchline main -> data :: --data
arrow data
node data :: Path("./data/emotions.csv")
arrow data
node data :: local filesystem
stop data

handoff main -> model :: --model
arrow model
node model :: Hugging Face model identifier
stop model
\end{dhkbranch}

그리고 내부에서는:

\begin{dhkcode}
sys.executable # 어떤 Python이 실행 중인지 알 수 있음
sys.path # Python이 import할 package를 어디서 찾는지 확인 할 수 있음
os.environ # 환경 변수에 접근할 수도 있음
\end{dhkcode}

이 정도가 되면 \textbf{CLI, Conda, VS Code, Jupyter, Python import, 파일 경로, URL, ML 학습 스크립트가 전부 하나의 시스템으로 보이기 시작한다.}

